\section{Weekly Summary}

\begin{tcolorbox}[colback=yellow!10,colframe=black!50,title=AI-Generated Content]
\noindent The summary below was written by Claude (Anthropic) from that week's regression outputs and series descriptions. It has not been reviewed or fact-checked by a human, and should be read as an automated, informal recap -- not analysis.
\end{tcolorbox}

Welcome back to the weekly parade of economic non-sequiturs, where a computer grabs two time series that have never met, forces them to sit next to each other, and we all pretend to be shocked by the results. This week we paired inflation with a discontinued interest rate, Swiss currency intervention with unfilled airplane orders, Asian air freight with fast-food job listings, and, in the grand finale, telephone-service prices with rich people's deposits. Truly, the questions humanity has always yearned to ask.

Let us begin with Monday's headliner: services inflation versus the tragically deceased 1-Week AMERIBOR rate. The raw regression came roaring in with an R-squared of 0.93, which if you didn't know better would look like the discovery of the century. But of course both of these things drift over time, so the real test is the detrended version, where the relationship politely deflates to an R-squared of about 0.23 -- still technically "significant," but more of a limp handshake than the crushing embrace the raw numbers promised. Classic trend mirage, only partially debunked. Then Tuesday happened, and the model tried to relate Swiss National Bank dollar purchases to unfilled transportation-equipment orders, only to discover the input series was apparently a flat line of zeros, producing an R-squared of roughly 0.0000000000000003 and enough "smallest eigenvalue is 0" warnings to make the software file a complaint. We regressed nothing on something and got, appropriately, nothing.

Midweek offered the reliable comfort of statistical failure. Asian air freight prices versus food-service job postings mustered a raw R-squared of 0.02 and a p-value that shrugged, and the detrended version somehow tried even less. Wednesday's loan sizes versus real gross domestic income were, frankly, indifferent to each other in both the raw and detrended runs -- a rare case of two series that agree they have nothing to say. And Thursday's top-0.1\%-loan-share versus Czech consumer pessimism achieved that beautiful "nope in both columns" symmetry, the kind of clean null result that reminds you the universe usually isn't hiding a secret.

But save your applause for Friday, because we actually found one worth pointing at. The Producer Price Index for residential local telephone service versus the deposit share held by the merely-very-rich (90th to 99th percentile) put up a modest raw R-squared of 0.12 -- and then, plot twist, the detrended version got stronger, climbing to about 0.41 with a wildly tiny p-value and even flipping to a negative relationship. That is genuinely the interesting case: a link that survives having its shared trend surgically removed, which is the closest this project ever gets to "huh, maybe." I have no earthly theory for why landline phone prices should move opposite to affluent deposit shares, and I want to be crystal clear that I am not proposing one. This is, as always, a machine blindly pairing random data with zero causal claim and no statistical rigor worth defending; spurious correlation is the house specialty here, not the exception. Friday just happened to produce a correlation that's spurious with slightly better production values. See you next week for more.
